\documentclass[12pt]{article}
\usepackage[utf8]{inputenc}
\usepackage[T1]{fontenc}
\usepackage{hyperref}
\title{Aluo Language Resources for K-12 Curriculum Development}
\author{Compiled by Claude Sonnet 5.5 (claude-sonnet-5-5) for Andr\'as Kornai}
\date{}
\begin{document}
\maketitle

\section{Introduction}
Aluo (ISO 639-3 yna; Glottocode aluo1235) is a Loloish (Sino-Tibetan) language of the Yi people. Wikipedia gives about 25,000 speakers (2007) in north Wuding, Luquan and Yuanmou counties (Yunnan) and Huili and Miyi counties (Sichuan); Joshua Project instead places most Aluo in Shizong and Luliang counties. Other names include Laka, Gan Yi, Yala, Lila and Niluo. Sources disagree on the branch (Wikipedia: Northern Loloish; Asia Harvest: Southeastern Yi) and on the endangerment level. It should not be confused with Naluo Yi (ylo), Naruo or Asuo.

\section{Primary Educational Resources}
No primers, textbooks or school programmes for Aluo were found in search results. The language is not known to be taught in schools.

\section{Linguistic and Reference Materials}
\begin{itemize}
\item Wikipedia, Aluo language: \url{https://en.wikipedia.org/wiki/Aluo_language}
\item Wikipedia, Loloish languages: \url{https://en.wikipedia.org/wiki/Loloish_languages}
\end{itemize}

\section{Religious and Community Resources}
\begin{itemize}
\item Joshua Project, Aluo in China: \url{https://joshuaproject.net/people_groups/18391/CH}
\item Asia Harvest, Aluo people group profile (PDF): \url{https://people-groups.asiaharvest.org/people-group-profiles/yunnan/aluo.pdf}
\end{itemize}
No scripture translations were identified.

\section{Digital and Audio Resources}
None found.

\section{Curriculum Development Context}
Only demographic and encyclopedic information was found. A curriculum would need new documentation and community involvement. The conflicting location and classification claims above should be settled against primary linguistic sources before any planning.

\end{document}